\subsection{Adaptive precision from cumulative-counter queries}
\label{sec:counter-precision-method}
Consider a six-hour block with nonnegative hourly increments
$z=(z_1,\ldots,z_6)$ and retained cumulative registers $C_0,C_6$,
all expressed in the same energy units.
Writing $E=C_6-C_0$, the initial feasible set is
$F_E=\{z\geq0:\sum_{h=1}^6z_h=E\}$.
For two frozen forecast vectors $a,b$, their block MSE contrast is
\[
 D(z)=d+c^\top z,\qquad
 d=\frac16\sum_{h=1}^6(a_h^2-b_h^2),\quad
 c_h=\frac{2(b_h-a_h)}6.
\]
The permitted query at boundary $k\in\{1,\ldots,5\}$ returns
$C_k$, equivalently the prefix sum $s_k=C_k-C_0$.
The budget is at most two distinct internal registers per pair and block.
This interface retrieves cumulative registers. The coordinate-increment
recovery controls are documented separately in Online Resource~1, S2--S4.
The register queries are exact linear measurements covered by the general
information model. Adaptive scalar certification also appears in monotone
integral estimation \citep[Section~4.1]{gaillard2023}; here the target is an
affine contrast of interval increments under a fixed hard budget.

For integer boundaries $0\leq p<q\leq6$, define
$r(p,q)=\max_{p<h\leq q}c_h-\min_{p<h\leq q}c_h$.
With queried boundaries $Q$, let $0=q_0<\cdots<q_m=6$ include
the outer boundaries and $Q$, and put $s_0=0$, $s_6=E$.
Nonnegativity and the segment totals give the sharp width
\begin{equation}
 W_Q(z)=\sum_{\ell=1}^{m}
 (s_{q_\ell}-s_{q_{\ell-1}})r(q_{\ell-1},q_\ell).
 \label{eq:counter-width}
\end{equation}
Both endpoints are attainable by placing each segment's mass at a
minimum- or maximum-coefficient hour. The initial width is $W_0=Er(0,6)$.
Tasks with $W_0=0$ are already precise and require no queries.

\begin{proposition}[Two best-side splits for an arbitrary block]\label{prop:counter-two-query}
Let $n\geq3$, $E>0$, $x\geq0$, $\sum_{h=1}^n x_h=E$, and let the target be
$d+\sum_hc_hx_h$. Each query returns one internal prefix sum. Write
$r(p,q)=\max_{p<h\leq q}c_h-\min_{p<h\leq q}c_h$.
For first boundary $k$, put $\ell_k=r(0,k)$ and $r_k=r(k,n)$.
Define each of the following quantities whenever its corresponding
side has an eligible second boundary:
\[
 \alpha_k=\min_{1\leq j<k}\max\{r(0,j),r(j,k)\},\qquad
 \beta_k=\min_{k<j<n}\max\{r(k,j),r(j,n)\}.
\]
Choose a minimizing boundary on each side, breaking ties by smallest index.
When both sides are nonempty, for normalized first answer $t=s_k/E$,
the optimal conditional worst width is
\[
 E\min\{\alpha_k t+r_k(1-t),\ \ell_k t+\beta_k(1-t)\}.
\]
Thus only one best boundary on each side is needed, independently of $t$.
If $\Delta_k=\ell_k+r_k-\alpha_k-\beta_k>0$, the switch is
\[
 t_k^\star=\frac{r_k-\beta_k}{\Delta_k}\in[0,1],
\]
and the worst width for this first boundary is
\[
 A_k=E\max\left\{\alpha_k,\beta_k,
          \frac{\ell_kr_k-\alpha_k\beta_k}{\Delta_k}\right\}.
\]
Use the best right split below the switch and the best left split above
it; either is valid at equality. If $\Delta_k=0$, the two lines coincide
and $A_k=E\max\{\ell_k,r_k\}$. At $k=1$ only the right split exists and
$A_1=E\beta_1$; at $k=n-1$ only the left split exists and
$A_{n-1}=E\alpha_{n-1}$. The optimal deterministic two-query guarantee is
\begin{equation} A=\min_{1\leq k<n} A_k.\label{eq:counter-adaptive-value}\end{equation}
\end{proposition}
\begin{proof}
All second queries left of $k$ have the common right-side contribution
$r_k(1-t)$; minimizing their remaining term $t$ times a split range
selects the same minimizing boundary for every $t\geq0$. The right-side
argument is identical. Moreover $\alpha_k\leq\ell_k$ and
$\beta_k\leq r_k$. Hence $\Delta_k\geq0$, and when it is positive the
displayed intersection belongs to $[0,1]$. The maximum of the lower
envelope occurs at an endpoint or that intersection. Its endpoint
values are $\beta_k,\alpha_k$, and its intersection value is the displayed
quotient. If $\Delta_k=0$, both nonnegative differences vanish, so the
lines coincide. At an outer eligible boundary the single-hour side has
zero coefficient range, giving the stated one-line values.
\end{proof}

All contiguous ranges can be built in $O(n^2)$ arithmetic operations
and stored in $O(n^2)$ space. Scanning the candidates for every
$\alpha_k,\beta_k$ costs another $O(n^2)$ operations; the closed forms
then cost $O(n)$. This is an arithmetic-operation bound for the exact-total,
uncapped two-prefix problem, not a bit-complexity or noisy-polytope claim.
For $n=2$, the single internal register determines both increments;
for $n=1$ the total already determines the only increment.


\paragraph{An upper certificate for randomized-batch precision.}
Returning to the six-hour case, there are ten two-boundary batches.
Let $z^{(h)}=Ee_h$ be the six feasible
single-hour mass profiles, and define
$M_{hQ}(\tau)=\mathbf1\{W_Q(z^{(h)})\leq\tau\}$.
The finite-profile program
\begin{equation}
 U(\tau)=\min_{w\geq0,\,\sum_hw_h=1,\,u}u
 \quad\text{subject to}\quad
 \sum_hw_hM_{hQ}(\tau)\leq u\quad\text{for every batch }Q
 \label{eq:counter-precision-upper}
\end{equation}
upper-bounds the worst-case precision success of \emph{every} randomized
batch policy. Indeed, its minimum success over $F_E$ is at most its success
averaged over these profiles under $w$, which is at most
$\max_Q\sum_hw_hM_{hQ}$. This remains valid for batches with fewer queries:
adding an unused register cannot increase width. The exported $w$ therefore
provides a checkable upper certificate. The finite profile set need not
contain a worst completion for an indicator objective, so $U(\tau)$ is not
asserted to equal the full precision optimum in general. Below, write
$S_B^\star=S_\tau^{\mathrm{batch}}$ and
$S_A=S_\tau^{\mathrm{adaptive}}$ for the precision-success values defined
in \eqref{eq:risk_precision_success}, at this fixed instance and threshold.
If $A\leq\tau$ and $U(\tau)<1$, the adaptive policy succeeds for every
completion, whereas no randomized batch can do so. In floating-point
classification, batch widths within $10^{-9}W_0$ above the threshold count
as success, and adaptive certification requires a margin of
$10^{-9}W_0$; exact constructed boundary cases are treated analytically.

\paragraph{An exact counter example.}
Take $E=1$ and $c=(0,0,1,1,0,1)$, so $W_0=1$.
Query boundary $4$ first, obtaining $t=s_4$.
Batch $\{2,4\}$ has width $1-t$, while batch $\{4,5\}$ has width $t$.
Querying $5$ when $t<1/2$ and $2$ otherwise thus guarantees
width $\min\{t,1-t\}\leq1/2$.
Conversely, at $z=e_5$ only batch $\{4,5\}$ attains width below $1$;
that batch has width $1$ at $z=e_2$. No batch succeeds at both profiles
for any threshold $\tau<1$. Weights $w_2=w_5=1/2$ in
\eqref{eq:counter-precision-upper} give $U\leq1/2$.
An equal mixture of $\{2,4\}$ and $\{4,5\}$ succeeds with probability
at least $1/2$ on every completion when $\tau\geq1/2$.
Consequently, for every $\tau\in[1/2,1)$ the optimal adaptive worst-case
success is $1$ and the optimal randomized-batch worst-case success is
$1/2$.
